\clearpage
\section{Canonical rule selection and recursive runtime safety}
\label{sec:copying-selection}
Every contextual rewrite of the concrete combined machine now has a unique
\emph{recognized rule choice}. Independently chosen decompositions of one
source graph agree on the operation family, ordered arithmetic operands,
atom-comparison result and copy-head primitive. More precisely, arbitrary
copy matches at the same wire select the same payload head \emph{edge}; read
and write matches select the same store edge. These are properties of the
actual host graph, with no global acyclicity assumption. The proof permits
different matched contraction trees, different local interfaces and different
representation allowances. It does not assume an isomorphism between the
chosen matches.

This completes the source-selection part of matching independence. It does
\emph{not} yet identify the graphs produced by independently chosen rewrites:
in particular, the residual tree, copied payload, distributors and surviving
host must still be compared by a boundary-respecting recursive isomorphism.
The theorem records a copy's head occurrence, not a new proof that the entire
copied box forest is unique. These remaining successor-graph obligations are
retained in Appendix~B.1, along with refocusing and the specimen laws.

A second result proves recursive label safety through every complete finite
run. The source translation alphabet excludes stores; the runtime alphabet
adds stores because reference allocation creates them. Neither hidden focus
edges nor unfilled holes can appear, even in duplicated or opened box bodies.
The original rewrite relation, stability and copyability premises, proposition
terms, and Figure~16 translations are unchanged.

The dependency chain is: incidence-preserving occurrences; canonical heads
and ordered arguments; deterministic term-sharing and atom-store resolution;
unique rule recognition; then preservation of recursive labels through all
local rules, their context closure and finite runs. These proofs supply
premises needed by the paper's copying-machine determinism argument in
Section~10.3, rather than replacing that argument with an assumed interface.

\Needspace{12\baselineskip}
\subsection{Occurrences inside arbitrary hosts}

\termdef{match-embedding}
An occurrence $f:G\hookrightarrow H$ consists of a map of ports and an injective
map of edge occurrences. Mapping the ordered source and target lists of an edge
must give exactly the corresponding host lists. A primitive label $p$ in $G$
must remain $p$ in $H$, interpreted at the host's depth allowance. No global
injectivity of ports is assumed: attachment may identify boundary ports.
Edge injectivity is essential for preserving distinct store identities.
This interface describes an occurrence; it is not an assertion that an arbitrary
non-injective variable renaming preserves program meaning. It also does not
compare the contents of arbitrary box labels, so later uses are restricted to
primitive recognition and incidence tracing.

\lean{Copying/MatchEmbedding}{10}{16}{MatchEmbedding}

The identity occurrence uses identity maps. Composition composes both maps;
edge injectivity composes, primitive-label preservation is used twice, and
mapping an incidence list twice equals mapping the composite function.
A previously proved search transport supplies an occurrence by forgetting
its inverse and boundary classification. A recursive graph isomorphism gives
such a search transport, hence an occurrence.

\lean{Copying/MatchEmbedding}{20}{26}{MatchEmbedding.refl}

\lean{Copying/MatchEmbedding}{28}{35}{MatchEmbedding.comp}

\lean{Copying/MatchEmbedding}{37}{43}{MatchEmbedding.ofTransport}

\lean{Copying/MatchEmbedding}{45}{46}{MatchEmbedding.ofDeepIso}

Shallow plugging maps inserted edges to the right summand and ports by the
actual attachment map. Simultaneous expansion maps a filler edge $a$ to the
dependent pair $(e,a)$ for its slot $e$. Both edge maps are injective; their
incidence and primitive-label laws are the definitions of plugging and
expansion. The left and right occurrences in serial assembly use its existing
edge and port maps; their edge images are the two distinct slot tags. These
constructions require no safety predicate on the erased slot label.

\lean{Copying/MatchEmbedding}{48}{56}{MatchEmbedding.plug}

\lean{Copying/MatchEmbedding}{58}{66}{MatchEmbedding.expand}

\lean{Copying/MatchEmbedding}{68}{74}{MatchEmbedding.serialLeft}

\lean{Copying/MatchEmbedding}{76}{84}{MatchEmbedding.serialRight}

Membership in a source or target list survives mapping because membership
survives list mapping. An adjacency witness is a shared port in a target and
a source; mapping that witness preserves adjacency. These facts deliberately
use forward incidence preservation, not a false assumption that the attachment
map is globally injective on ports.

\lean{Copying/MatchEmbedding}{86}{88}{MatchEmbedding.source_mem}

\lean{Copying/MatchEmbedding}{90}{92}{MatchEmbedding.target_mem}

\lean{Copying/MatchEmbedding}{94}{96}{MatchEmbedding.adjacent}

\Needspace{12\baselineskip}
\subsection{Canonical active heads and ordered arguments}

\termdef{matched-head}
$\mathsf{HeadAt}(G,v,e,p)$ says that edge $e$ has primitive label $p$ and consumes
wire $v$. If two heads consume the same wire, source-incidence uniqueness
forces their edges to be equal. Their labels are then equal, and injectivity
of primitive labelling forces the primitives to be equal. Transporting a head
uses primitive preservation and source membership preservation.

\lean{Copying/MatchEmbedding}{99}{100}{MatchEmbedding.HeadAt}

\lean{Copying/MatchEmbedding}{102}{106}{MatchEmbedding.head_unique}

\lean{Copying/MatchEmbedding}{108}{109}{MatchEmbedding.head}

\termdef{matched-argument}
$\mathsf{Argument}(G,e,i,a)$ says that the $i$th target of edge $e$ exists and
is consumed by edge $a$. Two witnesses use the same port because the ordered
list has one value at index $i$; source-incidence uniqueness then equates
the consuming edges. Mapping a list preserves the value at an existing
index, so occurrences preserve this ordered argument relation. This rules
out changing operand order while presenting the same host match.

\lean{Copying/MatchEmbedding}{112}{113}{MatchEmbedding.Argument}

\lean{Copying/MatchEmbedding}{115}{119}{MatchEmbedding.argument_unique}

\lean{Copying/MatchEmbedding}{121}{125}{MatchEmbedding.argument}

\Needspace{12\baselineskip}
\subsection{Term-sharing resolution}

\termdef{copy-resolution}
The forward labels for a copy trace are term contraction and term weakening.
A payload head is an instance or operation primitive. A hop follows an
adjacency whose departing edge has a forward label. Each forward edge has
exactly one term target, by its signature. Consequently two hops from that
edge agree on the next edge: their shared ports coincide, and source
uniqueness identifies the next consumer. A payload head is terminal for this
relation because neither an instance nor an operation is a sharing edge.
The weakening case is admitted in intermediate tree reasoning, although a
weakening cannot consume an externally selected input.

\lean{Copying/CopyMatching}{12}{13}{CopyTrace.ForwardLabel}

\lean{Copying/CopyMatching}{16}{16}{CopyTrace.PayloadHead}

\lean{Copying/CopyMatching}{18}{19}{CopyTrace.Hop}

\lean{Copying/CopyMatching}{21}{22}{CopyTrace.forward_output}

\lean{Copying/CopyMatching}{24}{26}{CopyTrace.hop_functional}

\lean{Copying/CopyMatching}{28}{34}{CopyTrace.head_terminal}

A resolution of wire $v$ to payload edge $e$, primitive $p$ and count $n$
chooses its first consuming edge and an exact $n$-hop trace ending at a
$p$-labelled payload head. Two successful resolutions first agree on their
starting edge by source uniqueness. The functional-terminal-trace theorem
then gives the same endpoint and exact length. Primitive-label injectivity
also gives the same $p$. A surrounding cycle causes no problem: the theorem
compares successful traces and does not claim every host trace terminates.

\lean{Copying/CopyMatching}{37}{38}{CopyTrace.Resolves}

\lean{Copying/CopyMatching}{41}{49}{CopyTrace.resolves_unique}

An occurrence maps sharing labels and adjacency, hence each hop, and maps
an entire exact trace without changing its length. Combining this transport
with uniqueness in the larger host proves that attachment cannot redirect an
already resolved request or supply a different successful hop count. The
auxiliary last-hop lemma appends one hop to a trace and increments its length;
its proof is induction on the trace.

\lean{Copying/CopyMatching}{51}{56}{CopyTrace.hop_map}

\lean{Copying/CopyMatching}{58}{62}{CopyTrace.resolves_map}

\lean{Copying/CopyMatching}{65}{68}{CopyTrace.resolves_map_exact}

\lean{Copying/CopyMatching}{71}{75}{CopyTrace.snoc}

\Needspace{12\baselineskip}
\subsection{Every actual copy match resolves to its payload}

A copyable graph has one term input, which source coverage assigns to an
edge. Its proved outer-label property permits only instances, operations or
boxes. Typing excludes a box at this input, since a box consumes a thunk wire.
Thus the edge is an instance/operation head. No condition on the behaviour or
stability of a lambda body is added to copyability.

\lean{Copying/CopyMatching}{80}{100}{Copyable.head_exists}

In the literal copying redex, the tree's output and the payload's input are
the same serial-assembly wire. Each tree adjacency becomes a sharing hop
through the left occurrence. The root targets that shared wire, so it has one
last hop to the payload head in the right occurrence. The arbitrary-tree
termination theorem supplies, for every selected input, a trace from its
consumer to the root. Mapping that trace and appending the last hop produces
resolution to the attached payload, with count $m+1$ for a tree trace of
length $m$.

\lean{Copying/CopyMatching}{105}{108}{CopyRule.tree_payload_wire}

\lean{Copying/CopyMatching}{111}{118}{CopyRule.before_tree_hop}

\lean{Copying/CopyMatching}{121}{131}{CopyRule.before_root_hop}

\lean{Copying/CopyMatching}{134}{144}{CopyRule.before_resolves}

The consumer of a tree input must be a contraction: a weakening's source
list is empty. This proves the active head required for copy recognition.
If a particular payload head was specified in advance, head uniqueness
identifies it with the head chosen by the existence proof, giving the same
resolution theorem for that specified edge.

\lean{Copying/CopyMatching}{147}{161}{CopyRule.before_active_head}

\lean{Copying/CopyMatching}{163}{170}{CopyRule.before_resolves_head}

\termdef{copy-match-head-independence}
Now place two independently chosen redexes into one host, with the same image
of their selected input. Their trees may have different input counts and need
not be isomorphic. Their payloads may be presented at different allowances.
Resolve each selected input in its own redex, map both traces into the host,
and apply host resolution uniqueness. The images of the payload heads are
the same edge and have the same primitive. This is the precise copy-selection
independence established here; it does not equate the two rewritten graphs.

\lean{Copying/CopyMatching}{174}{189}{CopyRule.match_payload_unique}

\Needspace{12\baselineskip}
\subsection{Read and write destinations}

Atom hops and successful store-resolution traces map through the same
occurrence interface. Unlike copy traces, their forward labels also include
instance edges, and their terminal edge is a store. The proofs map the exact
trace and preserve the store label; previously proved store-resolution
uniqueness supplies its canonical endpoint and count.

\lean{Copying/StoreMatching}{9}{15}{StoreTrace.hop_map}

\lean{Copying/StoreMatching}{17}{19}{StoreTrace.resolves_map}

\termdef{store-operand-match}
A resolved store operand at wire $v$ specifies an operation $o$, an argument
index $i$, a store edge $s$ and a hop count. It witnesses the operation head,
the edge consuming its $i$th argument, and a store-resolution trace from that
edge. Two witnesses agree on the head, then on the ordered argument edge,
then on the store and count. Occurrences preserve all three witnesses.

\lean{Copying/StoreMatching}{22}{23}{StoreTrace.Operand}

\lean{Copying/StoreMatching}{26}{34}{StoreTrace.operand_unique}

\lean{Copying/StoreMatching}{36}{40}{StoreTrace.operand_map}

For each concrete dereference redex, its reader's argument is the displayed
instance and its already proved trace resolves to the retained store. For
assignment, argument zero is the displayed instance resolving to the
destination store; the new stable value remains argument one. These witnesses
are built from the actual local incidence lists and existing tree-resolution
proofs. Given two independent embeddings of read redexes, or two of write
redexes, at the same host token, transport the operand witnesses and apply
uniqueness. Both matches select the same store occurrence. This result does
not yet compare the entire written-value graphs or the successor heaps.

\lean{Copying/StoreMatching}{45}{49}{Dereference.before_operand}

\lean{Copying/StoreMatching}{52}{61}{Dereference.match_store_unique}

\lean{Copying/StoreMatching}{66}{70}{Assignment.before_operand}

\lean{Copying/StoreMatching}{73}{83}{Assignment.match_store_unique}

\Needspace{12\baselineskip}
\subsection{Reference-equality branch independence}

\termdef{atom-match-decision}
A comparison decision $b$ at $v$ consists of resolved stores $s,t$ for the
ordered arguments zero and one, together with
\[
       s=t\quad\Longleftrightarrow\quad b=\mathsf{true}.
\]
Two decisions agree on each store by operand uniqueness, hence agree on $b$.
Mapping a decision uses edge injectivity in the reverse implication
$f(s)=f(t)\Rightarrow s=t$. This is why mere forward incidence preservation
would not suffice to exclude a false aliasing match.

\lean{Copying/StoreMatching}{88}{90}{AtomEquality.Decision}

\lean{Copying/StoreMatching}{93}{101}{AtomEquality.decision_unique}

\lean{Copying/StoreMatching}{104}{108}{AtomEquality.decision_map}

The one-store local rule has decision true: both existing atom traces end
at its same store. The two-store rule has decision false: the store edges
have distinct constructor tags in the actual frame. The operation head and
both ordered argument incidences are checked directly. If both redex shapes
could embed at the same host token, transported decisions would force true
to equal false. Thus the branches are disjoint under independently chosen
contexts and contraction trees, not just under a renaming of one match.

\lean{Copying/StoreMatching}{112}{120}{AtomEquality.Same.before_decision}

\lean{Copying/StoreMatching}{124}{135}{AtomEquality.Different.before_decision}

\lean{Copying/StoreMatching}{139}{148}{AtomEquality.same_different_disjoint}

\Needspace{12\baselineskip}
\subsection{Arithmetic operand independence}

\termdef{numeric-operand-match}
A number argument witnesses the edge consuming argument $i$ and requires its
label to be the actual natural-number constant $n$. The argument edge is
unique; equality of its primitive labels implies equality of the numbers.
Occurrences preserve this witness and its index.

\lean{Copying/ArithmeticMatching}{9}{10}{Arithmetic.NumberArgument}

\lean{Copying/ArithmeticMatching}{12}{19}{Arithmetic.numberArgument_unique}

\lean{Copying/ArithmeticMatching}{21}{24}{Arithmetic.numberArgument_map}

The local unary head is numeric negation and its argument is the specified
number edge. The binary head is the specified addition/subtraction operation;
its two number edges occur at indices zero and one respectively. All these
facts are checked from the original serial/tensor incidence lists, with no
evaluation or term simplification. The operation-code-to-symbol function is
injective by its two cases. Therefore independently embedded binary redexes
at one token agree on the operation code, the left number and the right
number. The arithmetic convention remains $-n=0-n$ on natural numbers.

\lean{Copying/ArithmeticMatching}{26}{27}{Arithmetic.unary_head}

\lean{Copying/ArithmeticMatching}{29}{33}{Arithmetic.unary_operand}

\lean{Copying/ArithmeticMatching}{35}{37}{Arithmetic.binary_head}

\lean{Copying/ArithmeticMatching}{39}{44}{Arithmetic.binary_left_operand}

\lean{Copying/ArithmeticMatching}{46}{51}{Arithmetic.binary_right_operand}

\lean{Copying/ArithmeticMatching}{53}{54}{Arithmetic.BinaryCode.symbol_injective}

\lean{Copying/ArithmeticMatching}{57}{69}{Arithmetic.binary_match_unique}

\Needspace{12\baselineskip}
\subsection{Unique selection for the full contextual relation}

\termdef{combined-rule-choice}
A rule choice records allocation, read, assignment, beta, a Boolean comparison
result, a copy-head primitive, a unary numeric operand, or a binary operation
code with both numbers. Its head function gives the corresponding active
operation or term contraction. Comparison choices with different answers,
copy choices with different payload labels, and arithmetic choices with
different numbers deliberately share a head; incidence and trace facts must
distinguish those cases.

\lean{Copying/CombinedSelection}{10}{15}{Combined.Choice}

\lean{Copying/CombinedSelection}{17}{25}{Combined.Choice.head}

\termdef{combined-recognition}
Recognition of a choice means that the source wire has its required head,
plus the appropriate comparison decision, copy resolution, or ordered number
arguments. Allocation, read, assignment and beta need no additional data in
this particular record. Their finer store-selection results were proved
separately above. Recognition only inspects the source graph; it neither
changes the operational semantics nor asserts that a recognized choice has
an applicable rewrite. For example, it does not remove beta's stable-argument
premise.

\lean{Copying/CombinedSelection}{28}{34}{Combined.Recognizes}

Recognition transports through occurrences by the head, argument and trace
transport lemmas. For uniqueness, head uniqueness first excludes differing
operation families. The remaining cases use comparison-decision uniqueness,
copy-resolution uniqueness, or number-argument uniqueness. Binary-symbol
injectivity handles addition versus subtraction. These exhaust the choice
constructors.

\lean{Copying/CombinedSelection}{37}{46}{Combined.recognizes_map}

\lean{Copying/CombinedSelection}{49}{69}{Combined.recognizes_unique}

Every local rewrite has a recognition witness. The fixed operation heads
are read directly from the local frames. Equality uses the true/false decision
lemmas, copying uses the arbitrary-tree head and resolution theorems, and
arithmetic uses the two ordered operand witnesses. The original local-rule
proof supplies its existing stability or copyability premise.

Induction on an actual contextual rewrite then lifts recognition through
recursive renaming, shallow plugging and simultaneous expansion. The unchanged
fillers need not be inspected for this source property. Combining existence
with uniqueness gives the milestone theorem:
\[
 G@v\longrightarrow H@w\quad\Longrightarrow\quad
 \exists!c,\ \mathsf{Recognizes}(G,v,c).
\]
This statement applies to every derivation of the existing rewrite relation.
Uniqueness of $c$ must not be confused with uniqueness of $(H,w)$ up to
representation; the latter still needs the successor-coherence argument.

\lean{Copying/CombinedSelection}{72}{90}{Combined.local_recognizes}

\lean{Copying/CombinedSelection}{93}{99}{Combined.rewrite_recognizes}

\lean{Copying/CombinedSelection}{103}{106}{Combined.rewrite_unique_choice}

Every choice head is active or a term contraction, by cases on the choice.
Stable graphs contain only instances and passive operation heads, so they
support no contextual rewrite at any wire. More generally, an embedded
stable payload cannot be rewritten at the image of its input, even if its
host is cyclic or contains other redexes: source coverage supplies its passive
head, whose image conflicts with every recognized rewrite head. Finally, a
rewrite cannot be based at an external output, because its recognized head
would consume that output, contradicting graph incidence coverage.

\lean{Copying/CombinedSelection}{109}{111}{Combined.choice_head_active}

\lean{Copying/CombinedSelection}{114}{123}{Combined.stable_no_rewrite}

\lean{Copying/CombinedSelection}{126}{128}{Combined.rewrite_not_output}

\lean{Copying/CombinedSelection}{130}{142}{Combined.stable_occurrence_no_rewrite}

For a full search transport, apply occurrence transport in both directions;
the inverse port map cancels the forward map. This proves recognition
preservation and reflection. Instantiating with arbitrary finite padding
shows that the chosen operation, copy-head label, comparison answer and
numeric operands do not change with unused representation allowance.

\lean{Copying/CombinedSelection}{145}{152}{Combined.recognizes_transport_iff}

\lean{Copying/CombinedSelection}{155}{157}{Combined.recognizes_pad_iff}

\Needspace{12\baselineskip}
\subsection{Recursive label preservation through computation}

\termdef{rule-label-admissibility}
A primitive predicate $P$ admits the local rules when it includes every
operation, instance, store, weakening and contraction label. This is a
sufficient closure condition for the new labels introduced by computation.
It need not admit focus or hole labels. Safety is the existing recursive
predicate $\mathsf{AllLabels}(P,G)$, so it inspects every box body as well as
outer labels.

\lean{Copying/CombinedSafety}{10}{15}{RuleSafety.Admits}

Safety of a serial composite implies safety of each component by evaluating
the predicate on its corresponding edge occurrence and using label equality.
For beta, source safety additionally gives safety of the chosen lambda body:
inspect its actual box edge, unfold the recursive label predicate and undo
the provably equal interface cast. Inspecting the argument occurrence gives
its complete recursive safety independently of the separate stability premise.

\lean{Copying/CombinedSafety}{17}{21}{RuleSafety.serial_left}

\lean{Copying/CombinedSafety}{23}{27}{RuleSafety.serial_right}

\lean{Copying/CombinedSafety}{30}{34}{RuleSafety.application_body}

\lean{Copying/CombinedSafety}{37}{41}{RuleSafety.application_argument}

Local preservation is proved by cases on all nine rule constructors.
Allocation extracts safety of the old value and adds only admitted prefix
labels. Read and equality use their existing all-label lemmas. Assignment
extracts safety of the value occurrence and applies its after-graph lemma.
Beta extracts the body and argument predicates and uses preservation when
opening the box. Copying extracts payload safety from the serial source and
preserves it in both copies and the environment distributors. Arithmetic
returns an admitted numeric constant.

For context closure, recursive isomorphisms transport the predicate. Plugging
splits each target edge into a surviving host edge, certified by source safety,
or an inserted edge, certified by induction. Simultaneous expansion uses
induction at the rewritten slot and the explicit equality of all other
fillers. No safety assumption on an erased host slot is needed.

\lean{Copying/CombinedSafety}{44}{63}{RuleSafety.local_allLabels}

\lean{Copying/CombinedSafety}{66}{85}{RuleSafety.rewrite_allLabels}

\Needspace{12\baselineskip}
\subsection{Runtime alphabet and reachable-state guarantees}

\termdef{copying-runtime-alphabet}
The runtime predicate is the union of the original translation alphabet with
the store primitive. A runtime graph satisfies this predicate at every depth.
All five admissibility fields follow by cases on their constructor. Every
translated-language graph is a runtime graph: prove this by induction on the
box allowance, injecting primitive membership into the union and recursively
handling boxed bodies.

The predicates excluding every focus label and excluding every hole label
also admit all local rules, by constructor disjointness. Applying the generic
rewrite theorem gives preservation of runtime membership, focus freedom and
hole freedom separately. This proves a runtime invariant; it does not
incorrectly claim that allocation preserves the smaller translation alphabet.

\lean{Copying/CombinedSafety}{88}{88}{RuleSafety.RuntimePrimitive}

\lean{Copying/CombinedSafety}{90}{90}{RuleSafety.RuntimeGraph}

\lean{Copying/CombinedSafety}{92}{98}{RuleSafety.admits_runtime}

\lean{Copying/CombinedSafety}{101}{109}{RuleSafety.language_runtime}

\lean{Copying/CombinedSafety}{111}{112}{RuleSafety.admits_focusFree}

\lean{Copying/CombinedSafety}{114}{115}{RuleSafety.admits_holeFree}

\lean{Copying/CombinedSafety}{120}{122}{Combined.rewrite_runtime}

\lean{Copying/CombinedSafety}{124}{125}{Combined.rewrite_focusFree}

\lean{Copying/CombinedSafety}{127}{128}{Combined.rewrite_holeFree}

A combined search step leaves the graph unchanged; a rewrite step uses the
proved label-preservation theorem. Induction on exact finite machine runs
therefore preserves any admitted recursive predicate. For an unbounded run,
use the already proved exact bounded-run representation, then apply that
induction at the fixed allowance. A certified closed program supplies initial
focus freedom and hole freedom, giving both at every reachable state.
If its initial graph is additionally in the translation alphabet, the same
argument gives runtime-alphabet membership throughout. These conclusions
hold for every finite prefix, irrespective of whether the execution eventually
terminates.

\lean{Copying/CombinedSafety}{130}{134}{Combined.step_allLabels}

\lean{Copying/CombinedSafety}{136}{140}{Combined.run_allLabels}

\lean{Copying/CombinedSafety}{143}{148}{Combined.reachable_safe}

\lean{Copying/CombinedSafety}{151}{155}{Combined.reachable_runtime}

\Needspace{12\baselineskip}
\subsection{Regression consequences, audit and remaining effort}

The arithmetic regression excludes two matches whose unequal operands were
swapped: binary-match uniqueness would force the two numbers equal. An
allocation regression proves that its after-graph really leaves the
translation alphabet by inspecting the newly created store; another proves
that it nevertheless remains in the runtime alphabet. A numeric value
embedded in any host cannot support a rewrite at its input, by the stable
occurrence theorem.

\lean{Copying/SelectionExamples}{9}{15}{SelectionExamples.unequal_operands_not_swappable}

\lean{Copying/SelectionExamples}{18}{23}{SelectionExamples.allocation_not_translation_alphabet}

\lean{Copying/SelectionExamples}{26}{30}{SelectionExamples.allocation_runtime}

\lean{Copying/SelectionExamples}{33}{36}{SelectionExamples.embedded_number_no_rewrite}

For each unchanged Proposition~12.2 term, every finite run in every certified
closing context remains focus-free and hole-free. The quantification covers
arbitrary variable/atom environments and finite prefixes arising from
admissible divergent callbacks or repeated closure use. These are safety
invariants, not new equivalence or termination proofs for such callbacks;
the Figure~22 template argument is still required.

\lean{Copying/SelectionExamples}{40}{43}{SelectionExamples.privateCell_runs_safe}

\lean{Copying/SelectionExamples}{46}{49}{SelectionExamples.constantResult_runs_safe}

All 103 new declarations are included in the axiom audit. The complete
build passes 369 jobs. The three audits check 6 abstract, 477 graph and 1104
copying declarations and use only \texttt{propext}, \texttt{Classical.choice}
and \texttt{Quot.sound}. No unfinished proof or project-specific axiom is
introduced.

Appendix~B.1 removes source-selection independence and recursive label safety
from the unfinished obligations, while retaining successor-graph coherence,
refocusing, focused-state correspondence, specimen locality, semantic
renaming/substitution, all required template robustness proofs and final
assembly. None of the four broad remaining packages is claimed wholly closed.

The estimate for completion of Proposition~12.2 is now about 80\% of the
mathematical effort, with a broad plausible range of 62--90\% (26 October 2026).
The previous estimate was 78\%. The increase reflects canonical selection for
all rewrite families and complete-run recursive safety. This is a judgement
about proof difficulty, not a declaration count; successor coherence and the
eleven-template argument still carry substantial uncertainty.

